% ======================================================================
% DRAFT: Computational model section
% Assumes the preamble macros from the modeling notes (\ys, \us, \den, \normal,
% \argmax, bbm, tikz bayesnet). Experiment numbering assumed:
%   Experiment 1 = study 6 (generic / specific / baseline)
%   Experiment 2 = study 8 (graded mix of generics and specifics)
%   Experiment 3 = study 9 (feature-localized generics)
% Lines marked TODO need a fact or a citation filled in.
% ======================================================================

\section{Computational model}

Our account has three ingredients, each corresponding to one of the open
questions raised above. First, a learner represents a category not by a single
inductive potential but by a \emph{coherence function} over the space of
possible features, which says how likely each feature is to be kind-linked: true
of the category as such rather than of particular members. Second, the learner
updates this function by interpreting what a speaker said, reasoning about why
the speaker chose a generic or a specific statement. Third, the coherence
function is smooth over a representation of world knowledge about features, so
that evidence about one feature spreads to similar features. Below we describe
the generative model the learner is assumed to hold, how each ingredient is
formalized, the qualitative predictions that follow, and how we fit the model to
participants' judgments and test each ingredient by lesioning it.

\subsection{Overview}

Consider a learner who hears a set of utterances $\us = \{u_1, \ldots, u_n\}$
about a novel category (e.g., Zarpies). Each utterance $u_i$ mentions a feature
$f_i$ and is either generic (``Zarpies $f_i$'') or specific (``This Zarpie
$f_i$''). The learner is then asked about a novel feature $f'$ that was never
mentioned and judges its prevalence $p' \in [0, 1]$: what proportion of Zarpies
have $f'$?

We assume that the learner explains the utterances by appeal to a latent
variable for each feature, $z_i \in \{0, 1\}$, indicating whether $f_i$ is
kind-linked for this category. Kind-linked features are ones the category
supports generalizing (if being absent-minded is kind-linked for professors,
then a new professor is likely absent-minded); features that are not kind-linked
are idiosyncratic to individuals. The probability that each feature is
kind-linked is given by a coherence function $\phi$, which assigns every
possible feature a real-valued coherence $\phi(f)$, with
$P(z = 1 \mid \phi) = \sigma(\phi(f))$ and $\sigma(y) = 1/(1 + e^{-y})$. A
category of high inductive potential is one whose coherence function is high
across much of feature space; a minimal category is one whose coherence is low
everywhere.

The learner's inference then has two steps (Figure~\ref{fig:graphical-model}).
Utterances provide evidence about which mentioned features are kind-linked, and
so about the coherence function:
\begin{equation}
    P(\phi \mid \us) \propto P(\phi) \prod_{i=1}^{n} \sum_{z_i \in \{0,1\}}
        P(u_i \mid z_i)\, P(z_i \mid \phi).
    \label{eq:phi-posterior}
\end{equation}
The coherence function in turn predicts whether the novel feature is kind-linked
and hence how prevalent it is:
\begin{equation}
    P(p' \mid \us) = \int P(\phi \mid \us) \sum_{z' \in \{0,1\}}
        P(p' \mid z')\, P(z' \mid \phi)\, d\phi.
    \label{eq:prevalence-predictive}
\end{equation}
Three terms in these equations carry the theory. The prior $P(\phi)$ encodes
world knowledge about which features go together (Section~\ref{sec:feature-space}).
The utterance likelihood $P(u_i \mid z_i)$ encodes how statements are
interpreted, literally or pragmatically (Section~\ref{sec:utterances}). The
prevalence likelihood $P(p' \mid z')$ links the abstract notion of
kind-linkedness to the graded prevalence judgments people actually give
(Section~\ref{sec:linking}).

% (graphical model figure from the notes goes here: phi -> z_i -> u_i, phi -> z' -> p')

This structure is hierarchical in the sense of \citet{kempetal2007learningoverhypotheses}:
the coherence function is an overhypothesis that sits above individual features
and is shared by all of them. That is what lets statements about mentioned
features bear on unmentioned ones. It differs from earlier models of generic
language \citep{tesslergoodman2019languagegeneralization} in its target. Those
models explain how ``Zarpies eat flowers'' changes beliefs about the prevalence
of eating flowers. Ours explains how it changes beliefs about the prevalence of
every other feature.

\subsection{Feature space and the coherence function}
\label{sec:feature-space}

\paragraph{Representing features.} To give the learner world knowledge about how
features relate, we represent each feature by an embedding of its description
from a sentence-embedding model trained on large text corpora
(TODO: model name and citation; 384-dimensional), projected to two dimensions
(TODO: projection method). Features with similar meanings (e.g., ``like to
dance'' and ``like to give high-fives'') lie close together, and features from
different domains (dietary preferences, personality, physical traits) occupy
different regions. We write $\mathbf{x}_i \in \mathbb{R}^2$ for the embedding of
feature $f_i$ and $\mathbf{x}'$ for that of the novel feature. This is a
deliberately minimal stand-in for world knowledge. It captures topical
similarity between features but not, for example, causal relations among them
\citep{rehder2003causalmodel}. % TODO check key

\paragraph{A Gaussian process prior over coherence.} We place a Gaussian process
(GP) prior on the coherence function \citep{rasmussenwilliams2006gaussian}.
% TODO add bib
For any finite set of features, here the $n$ mentioned features and the novel
feature, the vector of their coherences
$\ys = (y', y_1, \ldots, y_n)$ with $y_i = \phi(f_i)$ is multivariate normal:
\begin{equation}
    \ys \sim \normal\big(\mu_0 \mathbf{1},\; \mathbf{K}_\theta\big), \qquad
    [\mathbf{K}_\theta]_{ij} = s^2 \exp\!\left(-\frac{\lVert \mathbf{x}_i - \mathbf{x}_j \rVert^2}{2\ell^2}\right).
    \label{eq:gp-prior}
\end{equation}
The prior mean $\mu_0$ sets how coherent the learner expects a new category to
be before hearing anything; the output scale $s$ sets how far coherence can move
from that default; and the length scale $\ell$ sets how far evidence about one
feature spreads through feature space.

The length scale is the central quantity for our purposes, because its two
limits correspond to two different theories of inductive potential. As
$\ell \to \infty$, all features are perfectly correlated: the coherence function
collapses to a single number for the whole category, and evidence about any
feature raises or lowers expectations about every other feature equally. This is
the classic notion of inductive potential as a property of the category alone.
For finite $\ell$, the coherence function varies over feature space: hearing that
Zarpies have some personality trait raises the coherence of other personality
traits more than that of physical traits. Estimating $\ell$ from people's
judgments, and comparing the fitted model with one where $\ell$ is fixed to be
much larger than the feature space (Section~\ref{sec:lesions}), therefore tests
whether people's inferences about inductive potential are feature-specific.

\subsection{Interpreting utterances}
\label{sec:utterances}

The utterance likelihood $P(u_i \mid y_i) = \sum_{z_i} P(u_i \mid z_i) P(z_i \mid y_i)$
connects what was said to coherence. We consider a literal and a pragmatic
version.

\paragraph{Literal semantics.} Under a simple truth-conditional semantics, a
generic statement ``Zarpies $f_i$'' is true only if $f_i$ is kind-linked, whereas
a specific statement ``This Zarpie $f_i$'' is true either way, since any
individual can have an idiosyncratic feature:
\begin{equation}
    \den{u_i}(z_i) = \begin{cases}
        \mathbbm{1}[z_i = 1] & \text{if } u_i \text{ is generic,} \\
        1 & \text{if } u_i \text{ is specific.}
    \end{cases}
\end{equation}
A literal learner treats each statement as a noisy report of its truth,
$P(u_i \mid z_i) = (1 - \delta)\den{u_i}(z_i) + \delta\,(1 - \den{u_i}(z_i))$,
with a small error rate $\delta = 0.05$. A literal learner thus becomes more
confident that a feature is kind-linked after a generic, but a specific
statement is equally likely under both values of $z_i$ and so carries no
information about coherence.

\paragraph{Pragmatic interpretation.} A pragmatic learner instead asks why the
speaker chose this utterance over its alternative, following the Rational Speech
Act framework \citep{goodmanfrank2016pragmaticlanguage}. The learner models a
speaker who knows whether $f_i$ is kind-linked and chooses between the generic
and the specific to convey that to a literal listener $L_0$:
\begin{align}
    L_0(z_i \mid u_i) &\propto \den{u_i}(z_i)\, P(z_i \mid y_i), \\
    S_1(u_i \mid z_i) &= (1 - \epsilon)\,
        \frac{\exp\{\beta\, L_0(z_i \mid u_i)\}}{\sum_{u} \exp\{\beta\, L_0(z_i \mid u)\}}
        + \frac{\epsilon}{2},
    \label{eq:speaker}
\end{align}
where $\beta > 0$ is the speaker's rationality, $\epsilon = 0.05$ allows for
occasional uninformative choices, and the literal listener's prior is the
learner's current coherence for that feature.
% NOTE: utility is L0 probability (as implemented), not log L0 as in standard RSA.
The pragmatic learner then uses $P(u_i \mid z_i) = S_1(u_i \mid z_i)$.

Because a speaker who knows a feature is kind-linked would usually have said the
more informative generic, hearing a specific statement is now evidence that the
feature is \emph{not} kind-linked. The strength of this inference grows with
$\beta$. Pragmatic interpretation is therefore what allows specific statements
to lower inferred coherence, rather than leaving it unchanged.

\subsection{From coherence to prevalence judgments}
\label{sec:linking}

The final component, $P(p' \mid z')$, maps the abstract variable
``kind-linked'' onto the 0--100 prevalence scale participants used. Two
properties of these judgments shape it. First, features differ in how prevalent
people expect them to be regardless of the category: ``like to dance'' receives
higher ratings than ``can snap with their toes'' in every condition. Second,
ratings pile up at the ends of the scale, with many responses of exactly 0 or
100.

We therefore give each test feature $j$ its own pair of response distributions,
one for when it is kind-linked and one for when it is not, and model each as a
zero-and-one-inflated Beta distribution \citep{ospinaferrari2010inflatedbeta}:
% TODO add bib
\begin{equation}
    P(p' \mid z' = c) = \begin{cases}
        \pi_{0,j}^{c} & p' = 0, \\
        \pi_{1,j}^{c} & p' = 1, \\
        (1 - \pi_{0,j}^{c} - \pi_{1,j}^{c})\,
            \mathrm{Beta}(p' \mid m_j^{c} \kappa_j^{c},\, (1 - m_j^{c})\kappa_j^{c}) & 0 < p' < 1,
    \end{cases}
    \qquad c \in \{\text{KL}, \neg\text{KL}\}.
\end{equation}
Marginalizing over $z'$, the predicted rating distribution for feature $j$ is a
mixture of these two distributions weighted by the feature's coherence:
\begin{equation}
    P(p'_j \mid y'_j) = \sigma(y'_j)\, P(p'_j \mid z' = \text{KL})
        + \big(1 - \sigma(y'_j)\big)\, P(p'_j \mid z' = \neg\text{KL}).
    \label{eq:linking}
\end{equation}
To keep the two components interpretable, we constrain the kind-linked
component to sit above the other: its Beta mean is at least as high
($m_j^{\text{KL}} \ge m_j^{\neg\text{KL}}$), it puts at least as much mass on
100 and no more mass on 0. We also require each Beta part to be single-peaked so
that pile-ups at the endpoints are captured by the point masses. (Details are in
Appendix~\ref{app:linking}.)

This division of labor matters for interpreting the model. The linking function
fixes, for each feature, the range of ratings it can receive, which is the
segment between the mean of the non-kind-linked component and the mean of the
kind-linked one. Coherence determines where along that segment the feature falls
in a given condition. All of the model's predictions about how language changes
judgments, across conditions and across features within a condition, therefore
come from the coherence function and not from the linking function.
% NOTE: with 8 free params per feature the linking function alone fits each
% feature's overall mean; the theory is tested by within-feature condition
% differences and by the model comparisons below, not by overall r.

\subsection{Model predictions}
\label{sec:predictions}

Together, these components produce four qualitative predictions, which
Experiments 1--3 test.
% TODO: simulation figure, e.g. a 1-D feature space with coherence fields after
% generics / specifics / mixed input at small vs large ell, literal vs pragmatic.

\begin{enumerate}
    \item \textbf{Generic statements raise the prevalence of unmentioned
    features.} Generics raise the coherence of the mentioned features. Because
    the GP ties coherences together, this raises the coherence, and hence the
    predicted prevalence, of the novel feature. This requires only literal
    semantics plus the hierarchical structure.

    \item \textbf{Specific statements lower the prevalence of unmentioned
    features.} Under pragmatic interpretation, a specific statement is evidence
    against kind-linkedness, so it pulls coherence below the prior. A literal
    learner predicts no change relative to hearing nothing.

    \item \textbf{Effects are graded with the mix of generic and specific
    statements.} Each utterance contributes its own likelihood term
    (Eq.~\ref{eq:phi-posterior}), so inferred coherence, and with it predicted
    prevalence, rises with the proportion of generic statements.

    \item \textbf{Effects spread to similar features.} With a finite length
    scale, generics about features in one region of feature space raise the
    coherence of nearby novel features more than distant ones. A
    distance-blind learner ($\ell \to \infty$) predicts the same increase for
    every novel feature.
\end{enumerate}

\subsection{Lesioned models}
\label{sec:lesions}

To test whether each component is needed, we compare the full model with
lesioned versions that each remove one ingredient (Table~\ref{tab:lesions}).

\begin{table}[ht]
\centering
\caption{The full model and lesioned alternatives.}
\label{tab:lesions}
\begin{tabular}{lll}
\toprule
Model & Change from full model & Tests \\
\midrule
Full & --- & --- \\
Distance-blind & $\ell$ fixed at 10 ($\gg$ feature-space diameter $\approx 0.8$) & Prediction 4 \\
Literal & Literal semantics, no speaker model (no $\beta$) & Prediction 2 \\
No transfer & Utterances removed; coherence stays at the prior & Predictions 1--3 \\
\bottomrule
\end{tabular}
\end{table}

The distance-blind model keeps the full apparatus but makes coherence a single
category-level quantity, which tests whether feature similarity matters beyond
category-level inductive potential. The literal model tests whether pragmatic
reasoning is needed; because Experiment 3 contained only generic statements, for
which literal and pragmatic interpretation make similar predictions, we compare
these two only in Experiments 1 and 2. The no-transfer model tests whether
language affects judgments about unmentioned features at all.

\subsection{Parameter estimation and model comparison}
\label{sec:inference}

\paragraph{Data model.} Participants in the same condition (or, in Experiment 2,
participants who heard the same assignment of generic and specific statements to
features) form a group $g$. We assume that all participants in a group share
the coherence function implied by what they heard, and that each participant's
rating of each test feature is an independent draw from
Eq.~\ref{eq:linking}. So the coherence of a feature is the probability that a
participant treats it as kind-linked. For each group, the latent coherences of
the $J$ test features and the $n_g$ mentioned features, $\ys_g$, receive the GP
prior in Eq.~\ref{eq:gp-prior}.

\paragraph{Parameters.} The model has four free parameters shared across all
groups in an experiment, $\theta = (\ell, \mu_0, s, \beta)$, plus the
linking-function parameters for each test feature. We place weakly informative
priors on $\theta$:
$\log \ell \sim \normal(\log 0.5, 1.5^2)$,
$\mu_0 \sim \normal(0, 1)$,
$\log s \sim \normal(\log 1.5, 1)$, and
$\log \beta \sim \normal(\log 3, 1)$.
% TODO report bounds of the VBMC box if any posterior presses against them
% (Exp 3 ZOIB fit: beta median 25.9, box edge 30).

\paragraph{Marginal likelihood over coherence.} For a given $\theta$, the
likelihood of a group's data requires integrating over the latent coherences:
\begin{equation}
    P(\mathbf{r}_g, \us_g \mid \theta) = \int
        \prod_{k, j} P(r_{kj} \mid y'_j)\,
        \prod_{i} P(u_i \mid y_i)\;
        \normal(\ys_g \mid \mu_0\mathbf{1}, \mathbf{K}_\theta)\, d\ys_g,
\end{equation}
where $r_{kj}$ is participant $k$'s rating of test feature $j$. We approximate
this integral with a Laplace approximation \citep{rasmussenwilliams2006gaussian}:
we find the posterior mode $\hat{\ys}_g$ by Newton's method and approximate
the posterior by a Gaussian with precision equal to the negative Hessian
$\mathbf{H}_g$ of the log joint at the mode, giving
\begin{equation}
    \log P(\mathbf{r}_g, \us_g \mid \theta) \approx
        \log \tilde{p}(\hat{\ys}_g) + \tfrac{D_g}{2}\log 2\pi
        - \tfrac{1}{2}\log\lvert \mathbf{H}_g \rvert,
\end{equation}
where $\tilde{p}$ is the integrand above and $D_g = J + n_g$. Groups are
conditionally independent given $\theta$, so the total log likelihood is the
sum over groups.

\paragraph{Linking parameters.} Rather than integrating over the 8 linking
parameters of each test feature, we set them to their (lightly penalized)
maximum-likelihood values at each $\theta$, given the coherences the model
predicts from the utterances alone (a profile likelihood). Because the full
model and every lesioned model profile the same linking function, this choice
does not favor any of them.

\paragraph{Posterior over parameters.} We infer the posterior over $\theta$
with Variational Bayesian Monte Carlo \citep[VBMC;][]{acerbi2018vbmc},
% TODO add bib
which is designed for expensive likelihoods such as ours and also returns an
estimate of the log model evidence $\log P(\text{data} \mid \text{model})$
(the evidence lower bound). We compare models by the difference in log evidence
(the log Bayes factor). Because the evidence integrates over $\theta$, a model
with more free parameters, such as the full model relative to the
distance-blind one, pays for its added flexibility automatically.
% TODO: for lesions that change the utterance likelihood (literal, no transfer),
% compare log P(ratings | utterances, model), not the joint evidence. See notes.

\paragraph{Model predictions.} To evaluate how well the model predicts
judgments, we compute its predictions from the utterances alone: for each group
we find the coherence mode given only the GP prior and the utterances
(no ratings), and report the expected rating under Eq.~\ref{eq:linking} at that
coherence, evaluated at the posterior median of $\theta$. Ratings therefore
inform the model only through the shared parameters $\theta$ and the linking
function, never through the group-specific coherences being predicted.

% ======================================================================
\appendix
\section{Model details}

\subsection{Laplace approximation}
% Newton with backtracking line search, initialized at mu_0; jitter 1e-6 on the
% Hessian, 1e-4 on K. Predictive coherence at a new location given the
% training-block mode (used for coherence-field figures):
% E[y(x*)] = mu_0 + k_*^T K^{-1} (\hat y_train - mu_0 1).

\subsection{Linking function}
\label{app:linking}
% Parameterization (m, kappa, pi0, pi1 per class), constraints
% a,b > 1; kappa <= 100; m in [0.02, 0.98]; m_KL >= m_notKL;
% pi1_KL >= pi1_notKL; pi0_KL <= pi0_notKL; Beta(1.5, 1.5) penalty on each pi.
% Fit: damped Newton from multiple data-driven and random starts per feature.

\subsection{VBMC settings}
% budget 400 evaluations, fixed seed, convergence diagnostics per fit.
